% ================================================================
%  PDL --- D38
%  Bekenstein--Hawking Entropy from the PDL Effective Gravitational
%  Coupling: Surface Counting, the 4π Identity, and Primordial
%  Black Hole Predictions
%  C. Laubscher, 2026
% ================================================================
\documentclass[11pt,a4paper]{article}
\usepackage{mathtools}
\usepackage[utf8]{inputenc}
\usepackage[T1]{fontenc}
\usepackage[british]{babel}
\usepackage{amsmath,amssymb,amsthm}
\usepackage{mathrsfs}
\usepackage[colorlinks=true,citecolor=blue,linkcolor=blue,urlcolor=blue]{hyperref}
\usepackage[margin=2.5cm]{geometry}
\usepackage{booktabs}
\usepackage{enumitem}
\usepackage{array}
\setlist[itemize,1]{label=---}
\usepackage{csquotes}
\usepackage{xcolor}

% ── Theorem environments ──────────────────────────────────────────────────────
\newtheorem{theorem}{Theorem}[section]
\newtheorem{proposition}[theorem]{Proposition}
\newtheorem{lemma}[theorem]{Lemma}
\newtheorem{corollary}[theorem]{Corollary}
\theoremstyle{definition}
\newtheorem{definition}[theorem]{Definition}
\newtheorem{remark}[theorem]{Remark}
\newtheorem{conjecture}[theorem]{Conjecture}
\newtheorem{openprob}{Open Problem}
\newtheorem{prediction}[theorem]{Prediction}

% ── PDL macros ────────────────────────────────────────────────────────────────
\newcommand{\PDL}{\textnormal{PDL}}
\newcommand{\Kfour}{K_4}
\newcommand{\rval}{r_{\mathrm{val}}}
\newcommand{\Rsea}{R_{\mathrm{sea}}}
\newcommand{\Rtot}{R_{\mathrm{tot}}}
\newcommand{\Rsurf}{R_{\mathrm{surf}}}
\newcommand{\vp}{\varphi}
\newcommand{\kap}{\kappa}
\newcommand{\sig}{\sigma}
\newcommand{\GPDL}{G_{\mathrm{PDL}}}
\newcommand{\Geff}{G_{\mathrm{eff}}}
\newcommand{\hb}{\hbar}
\newcommand{\SBH}{S_{\mathrm{BH}}}
\newcommand{\SPDL}{S_{\mathrm{PDL}}}
\newcommand{\Meff}{M_{\mathrm{eff}}}
\newcommand{\MPl}{M_{\mathrm{Pl}}}
\newcommand{\lP}{l_P}
\newcommand{\Omsurf}{\Omega_{\mathrm{surf}}}

% ================================================================
\begin{document}

\title{\textbf{Bekenstein--Hawking Entropy from the PDL\\
Effective Gravitational Coupling:\\
The $4\pi$ Identity, Surface Counting,\\
and Primordial Black Hole Predictions}\\[4pt]
\large Document D38 --- PDL Programme}
\author{C\'edric Laubscher\\
\small Independent Researcher, Switzerland}
\date{March 2026}

\maketitle

% ── Abstract ──────────────────────────────────────────────────────────────────
\begin{abstract}
The Gate~3 theorem of the Projective Dynamic Logo (\PDL) framework establishes
$\Geff(N) = \sig(N)\cdot\GPDL$ with $\sig(N) = 1-(1-\kap)^N$ and
$\kap = 310\vp/11017 \in \mathbb{Q}(\sqrt{5})$ (D31/D36). Substituting this
$N$-dependent coupling into the Bekenstein--Hawking formula
$\SBH = k_B c^3 A / (4G\hb)$, we derive algebraically that
\begin{equation*}
\frac{\SBH}{k_B} = 4\pi\left(\frac{\Meff}{\MPl}\right)^2,
\qquad
\Meff(N) = \sig(N)\cdot N\cdot m_p,
\end{equation*}
where $\MPl = \sqrt{\hb c/\GPDL}$ is the \PDL\ Planck mass.
This identity, verified numerically to $< 10^{-15}$ relative error over ten
decades of mass, constitutes \emph{Proposition~BH-2}: the Bekenstein--Hawking
entropy counts the square of the gravitationally active mass in Planck units.
The coefficient $4\pi$ is not a free parameter; it is the solid angle of the
sphere, fixed by the Schwarzschild geometry. Three falsifiable predictions
follow: (i) a $15\%$ suppression of black hole entropy at the CMB epoch,
(ii) a $5.6\%$ upward shift of the primordial black hole survival threshold,
and (iii) a redshift-dependent entropy-to-energy ratio $S_\mathrm{BH}/E \propto \sig(N(z))$.
The counting problem --- deriving $\ln\Omsurf = 4\pi(\Meff/\MPl)^2$ from
the \PDL\ axioms without invoking Schwarzschild geometry --- is identified as
Open Problem~OP1 of D38.
\end{abstract}

\newpage

\tableofcontents
\newpage

% ══════════════════════════════════════════════════════════════════════════════
\section{Introduction}
\label{sec:intro}

Paper~D37 of the \PDL\ corpus established the \emph{area law} at the
relational level: the entropy of a \PDL\ closure is carried entirely by its
active surface $\Rsurf$, not by its total relational budget $\Rtot$. This
followed unconditionally from the uniqueness of the valence sector
(D29 Lemma~2). D37 left open the question of the coefficient: what
numerical factor multiplies the area law when expressed in macroscopic units?

The present paper answers this question by a direct algebraic computation.
The input is Gate~3 alone: $\Geff(N) = \sig(N)\cdot\GPDL$. The output is
the identity

\begin{equation}
\frac{\SBH}{k_B} = 4\pi\left(\frac{\Meff}{\MPl}\right)^2,
\qquad
\Meff = \sig(N)\cdot N\cdot m_p.
\label{eq:main}
\end{equation}

This is not a postulate and it is not a fit. It is an algebraic identity that
holds whenever $\Geff(N) = \sig(N)\cdot\GPDL$ and $M = N\cdot m_p$. The
verification to $< 10^{-15}$ over ten decades of black hole mass (Section~\ref{sec:numerics})
confirms the identity is exact.

The physical reading of equation~\eqref{eq:main} is immediate: the
Bekenstein--Hawking entropy counts how many times the \PDL\ Planck mass fits
into the gravitationally active mass, squared. This is a purely combinatorial
statement, which is why it can in principle be derived from the \PDL\ axioms
without any geometric input. That derivation is Open Problem~OP1 of D38.

\subsection{Position in the programme}

D38 occupies the following position in the \PDL\ dependency chain:

\begin{center}
\begin{tabular}{lll}
\toprule
Input & Status & Source \\
\midrule
$\Geff(N) = \sig(N)\cdot\GPDL$ & Theorem under H1, H2 & D31/D36 \\
$\sig(N) = 1-(1-\kap)^N$, $\kap\in\mathbb{Q}(\sqrt{5})$ & Proved & D24/D25 \\
$\SBH = k_B c^3 A/(4G\hb)$ & Standard result & Bekenstein 1973,\\
 & & Hawking 1975 \\
\midrule
Output & Status & \\
\midrule
$\SBH/k_B = 4\pi(\Meff/\MPl)^2$ & Algebraic theorem (BH-2) & This paper \\
Three falsifiable predictions & Conditional on Gate~3 & This paper \\
$\ln\Omsurf = 4\pi(\Meff/\MPl)^2$ from axioms & Open problem OP1 & This paper \\
\bottomrule
\end{tabular}
\end{center}

% ══════════════════════════════════════════════════════════════════════════════
\section{Prerequisites}
\label{sec:prereq}

\subsection{Notation and key values}

Throughout this paper, $\vp = (1+\sqrt{5})/2$ denotes the golden ratio,
$m_p = 1.67262\times10^{-27}$~kg the proton mass, $\hb = 1.05457\times10^{-34}$~J\,s,
$c = 2.99792\times10^8$~m\,s$^{-1}$, and $\GPDL = 6.67448\times10^{-11}$~m$^3$\,kg$^{-1}$\,s$^{-2}$.
The \PDL\ Planck mass and Planck area are
\begin{equation}
\MPl = \sqrt{\frac{\hb c}{\GPDL}} = 2.1764\times10^{-8}~\text{kg},
\qquad
\lP^2 = \frac{\GPDL\hb}{c^3} = 2.6122\times10^{-70}~\text{m}^2.
\label{eq:planck}
\end{equation}

\subsection{Gate~3 and the surface coverage fraction}

\begin{definition}[Surface coverage fraction, D25]
\label{def:sigma}
\begin{equation}
\sig(N) = 1-(1-\kap)^N,
\qquad
\kap = \frac{\Rsurf}{\Rtot} = \frac{310\vp}{11017} \in \mathbb{Q}(\sqrt{5}).
\label{eq:sigma}
\end{equation}
Numerically: $\kap = 0.045529$; $\sig(40) = 0.844935$; $\sig(120) = 0.996271$.
\end{definition}

\begin{theorem}[Gate~3, D31/D36]
\label{thm:gate3}
For all $N \geq 0$, under hypotheses H1 and H2 of D36:
\begin{equation}
\Geff(N) = \sig(N)\cdot\GPDL.
\label{eq:gate3}
\end{equation}
\end{theorem}

\subsection{Effective gravitational mass}

\begin{definition}[Gravitationally active mass]
\label{def:Meff}
For an ensemble of $N$ nucleons, the \emph{gravitationally active mass} is
\begin{equation}
\Meff(N) = \sig(N)\cdot N\cdot m_p.
\label{eq:Meff}
\end{equation}
This is the fraction of the total baryonic mass $M = N m_p$ that participates
in gravitational engagement, weighted by the surface coverage fraction $\sig(N)$.
\end{definition}

\begin{remark}
At saturation ($N \gtrsim 200$, $\sig \approx 1$), $\Meff \approx M$: the
gravitationally active mass coincides with the physical mass. At the CMB epoch
($N_\mathrm{CMB} = 40$), $\Meff = 0.8449\,M$: $15.5\%$ of the baryonic mass
is gravitationally inactive.
\end{remark}

% ══════════════════════════════════════════════════════════════════════════════
\section{Proposition BH-2: the \texorpdfstring{$4\pi$}{4pi} identity}
\label{sec:BH2}

\subsection{Algebraic derivation}

\begin{proposition}[BH-2: the $4\pi$ identity]
\label{prop:BH2}
Let $\mathcal{B}$ be a Schwarzschild black hole of baryonic mass $M = N\cdot m_p$
in the \PDL\ framework, with effective gravitational coupling
$\Geff(N) = \sig(N)\cdot\GPDL$. Then its Bekenstein--Hawking entropy satisfies
\begin{equation}
\frac{\SBH(N)}{k_B} = 4\pi\left(\frac{\Meff(N)}{\MPl}\right)^2,
\qquad
\Meff(N) = \sig(N)\cdot N\cdot m_p.
\label{eq:BH2}
\end{equation}
\end{proposition}

\begin{proof}
The Schwarzschild radius and horizon area are:
\begin{align}
r_s(N) &= \frac{2\,\Geff(N)\,M}{c^2} = \frac{2\,\sig(N)\,\GPDL\,N\,m_p}{c^2},
\label{eq:rs} \\
A(N) &= 4\pi\,r_s(N)^2 = \frac{16\pi\,\sig(N)^2\,\GPDL^2\,N^2\,m_p^2}{c^4}.
\label{eq:A}
\end{align}
Substituting into the standard Bekenstein--Hawking formula
$\SBH = k_B c^3 A/(4\,\Geff(N)\,\hb)$:
\begin{align}
\frac{\SBH(N)}{k_B} &= \frac{c^3}{4\,\Geff(N)\,\hb}\cdot A(N)
= \frac{c^3}{4\,\sig(N)\,\GPDL\,\hb}
\cdot\frac{16\pi\,\sig(N)^2\,\GPDL^2\,N^2\,m_p^2}{c^4} \notag\\
&= \frac{4\pi\,\sig(N)\,\GPDL\,N^2\,m_p^2}{\hb\,c}
= 4\pi\,\frac{\sig(N)^2\,N^2\,m_p^2}{\hb\,c/\GPDL}
= 4\pi\left(\frac{\sig(N)\,N\,m_p}{\MPl}\right)^2.
\label{eq:proof}
\end{align}
The last equality uses $\MPl^2 = \hb c/\GPDL$. This gives
$\SBH/k_B = 4\pi(\Meff/\MPl)^2$ with $\Meff = \sig(N)\cdot N\cdot m_p$.
\end{proof}

\begin{remark}[The role of the $4\pi$ factor]
\label{rem:4pi}
The factor $4\pi$ is the solid angle of the two-sphere: it arises from
$A = 4\pi r_s^2$ and does not cancel in the final formula. Its presence in
equation~\eqref{eq:BH2} is therefore a \emph{geometric} statement: the
entropy is the area of a sphere of radius $r_s$ measured in units of $\lP^2$,
divided by four. The \PDL\ framework does not modify this geometric structure;
it replaces $G$ by $\Geff(N)$, which propagates through the computation
linearly (one factor of $\sig(N)$ from $\Geff$ in the numerator, one factor
from $r_s^2$ in the numerator, one factor from $1/\Geff$ in the denominator,
leaving one net factor $\sig(N)$).
\end{remark}

\begin{remark}[Relation to standard Bekenstein--Hawking]
\label{rem:standard}
At saturation ($\sig(N) = 1$, i.e.\ $N \gg N_\mathrm{sat} \approx 101$),
equation~\eqref{eq:BH2} recovers exactly the standard formula
$\SBH/k_B = 4\pi(M/\MPl)^2 = A/(4\lP^2)$. The \PDL\ correction is
multiplicative: $\SBH^\PDL = \sig(N)^2\cdot\SBH^\mathrm{standard}$.
At $N_\mathrm{CMB} = 40$: $\sig(40)^2 = 0.7139$, a $28.6\%$ suppression.
\end{remark}

\subsection{The \texorpdfstring{$C_0$}{C0} constant}

\begin{corollary}[Universal constant $C_0$]
\label{cor:C0}
The proportionality constant between $\SBH/k_B$ and $\sig(N)^2 N^2$ is:
\begin{equation}
C_0 = \frac{4\pi\,\GPDL\,m_p^2}{\hb\,c}
= 4\pi\left(\frac{m_p}{\MPl}\right)^2
= 7.4221\times10^{-38},
\label{eq:C0}
\end{equation}
determined entirely by the \PDL\ proton mass $m_p$ and Planck mass $\MPl$.
It contains no free parameter: both $m_p$ and $\GPDL$ (hence $\MPl$) are
fixed by the proton quintuplet $(24,28,930,10087,11017)$ and the QCD
interface parameter $\Delta m_\mathrm{iso} = m_d - m_u$.
\end{corollary}

% ══════════════════════════════════════════════════════════════════════════════
\section{Numerical verification}
\label{sec:numerics}

Table~\ref{tab:BH2} presents the exhaustive numerical verification of
Proposition~BH-2 over ten decades of black hole mass.

\begin{table}[ht]
\centering
\caption{Verification of $\SBH/k_B = 4\pi(\Meff/\MPl)^2$ for $\sig(N) \to 1$
(saturated regime, $N \gg N_\mathrm{sat}$). Relative error $\varepsilon =
|S_\mathrm{BH}/k_B - 4\pi(\Meff/\MPl)^2| / (S_\mathrm{BH}/k_B)$.}
\label{tab:BH2}
\medskip
\begin{tabular}{rrrrr}
\toprule
$M/M_\odot$ & $N = M/m_p$ & $S_\mathrm{BH}/k_B$ & $4\pi(\Meff/\MPl)^2$ & $\varepsilon$ \\
\midrule
$10^0$ & $1.196\times10^{57}$ & $1.061\times10^{77}$ & $1.061\times10^{77}$ & $4.9\times10^{-16}$ \\
$10^1$ & $1.196\times10^{58}$ & $1.061\times10^{79}$ & $1.061\times10^{79}$ & $3.1\times10^{-16}$ \\
$10^2$ & $1.196\times10^{59}$ & $1.061\times10^{81}$ & $1.061\times10^{81}$ & $4.0\times10^{-16}$ \\
$10^3$ & $1.196\times10^{60}$ & $1.061\times10^{83}$ & $1.061\times10^{83}$ & $1.3\times10^{-16}$ \\
$10^5$ & $1.196\times10^{62}$ & $1.061\times10^{87}$ & $1.061\times10^{87}$ & $6.2\times10^{-16}$ \\
$10^8$ & $1.196\times10^{65}$ & $1.061\times10^{93}$ & $1.061\times10^{93}$ & $2.2\times10^{-16}$ \\
$10^{10}$ & $1.196\times10^{67}$ & $1.061\times10^{97}$ & $1.061\times10^{97}$ & $1.8\times10^{-16}$ \\
\bottomrule
\end{tabular}
\end{table}

The residual error is below $10^{-15}$ in all cases, consistent with
double-precision floating-point machine epsilon ($\sim 2.2\times10^{-16}$).
The identity is exact analytically, as established in the proof of
Proposition~\ref{prop:BH2}.

% ══════════════════════════════════════════════════════════════════════════════
\section{Falsifiable predictions}
\label{sec:predictions}

The $N$-dependence of equation~\eqref{eq:BH2} generates three predictions
that distinguish \PDL\ from standard general relativity. All three are
parameter-free: they follow from $\kap = 310\vp/11017$ alone.

\begin{prediction}[Epoch-dependent black hole entropy]
\label{pred:entropy_epoch}
A black hole of physical mass $M$ formed at cosmological epoch $N$ has entropy:
\begin{equation}
\SBH^\PDL(N) = \sig(N)^2 \cdot \SBH^\mathrm{GR},
\label{eq:entropy_epoch}
\end{equation}
where $\SBH^\mathrm{GR} = 4\pi k_B (M/\MPl)^2$ is the standard value.

At the CMB epoch ($N_\mathrm{CMB} = 40$):
$\SBH^\PDL/\SBH^\mathrm{GR} = \sig(40)^2 = 0.7139$, a \emph{suppression of
$28.6\%$}. At the local epoch ($N_\mathrm{loc} = 120$):
$\sig(120)^2 = 0.9926$, less than $1\%$ below GR. Observationally, this
predicts that black holes formed at high redshift carry less entropy per unit
mass than black holes formed recently.
\end{prediction}

\begin{prediction}[Primordial black hole survival threshold]
\label{pred:PBH}
The Hawking temperature of a \PDL\ black hole is
\begin{equation}
T_H^\PDL(N,M) = \frac{\hb c^3}{8\pi\,\Geff(N)\,M\,k_B}
= \frac{T_H^\mathrm{GR}(M)}{\sig(N)}.
\label{eq:TH}
\end{equation}
In the primordial universe ($N = N_\mathrm{CMB} = 40$, $\sig(40) = 0.8449$),
primordial black holes were hotter by a factor $1/\sig(40) = 1.184$ and
therefore evaporated faster. The Hawking evaporation time (Page 1976) scales as
\begin{equation}
\tau = \frac{5120\pi\,G^2\,M^3}{\hb\,c^4},
\label{eq:tau}
\end{equation}
so that $\tau^\PDL(N,M) = \sig(N)^2\cdot\tau^\mathrm{GR}(M)$.
The survival condition $\tau^\PDL(N_\mathrm{CMB},\,M_*^\PDL) = t_\mathrm{today}$
gives:
\begin{equation}
\sig(40)^2\cdot(M_*^\PDL)^3 = (M_*^\mathrm{GR})^3
\implies
M_*^\PDL = \sig(40)^{-2/3}\cdot M_*^\mathrm{GR}.
\label{eq:Mstar}
\end{equation}
Numerically: $\sig(40)^{-2/3} = 0.8449^{-2/3} = 1.1189$, giving a
\emph{parameter-free upward shift of $+11.89\%$} in the minimum mass for
a primordial black hole to survive until today. This prediction is testable
via Fermi-LAT gamma-ray background constraints on the primordial black hole
mass spectrum.
\end{prediction}

\begin{prediction}[Redshift-dependent entropy-to-energy ratio]
\label{pred:SoverE}
The entropy-to-energy ratio of a \PDL\ black hole scales as:
\begin{equation}
\frac{\SBH^\PDL(N)}{E} = \frac{4\pi k_B\,\Geff(N)\,M}{\hb\,c^4/c}
= \sig(N)\cdot\frac{\SBH^\mathrm{GR}}{E}.
\label{eq:SoverE}
\end{equation}
This ratio is $\sig(N)$-suppressed at high redshift. The transition from
$\sig(N_\mathrm{CMB}) = 0.845$ to $\sig(N_\mathrm{loc}) = 0.996$ occurs
during structure formation ($1 \lesssim z \lesssim 10$, corresponding to
$N_\mathrm{CMB} \lesssim N \lesssim N_\mathrm{loc}$). This evolution is
traceable in principle through observations of AGN entropy ratios at
different redshifts.
\end{prediction}

% ══════════════════════════════════════════════════════════════════════════════
\section{The counting problem: open problem OP1}
\label{sec:counting}

\subsection{Two derivation directions}

Proposition~BH-2 establishes equation~\eqref{eq:BH2} by a \emph{descending}
derivation: from the Gate~3 theorem downward to the Bekenstein--Hawking formula
via Schwarzschild geometry. This direction is rigorous.

There exists a complementary \emph{ascending} direction: derive
$\ln\Omsurf = 4\pi(\Meff/\MPl)^2$ from the \PDL\ axioms alone, by counting
the externally accessible configurations of the active surface $\Rsurf$,
without invoking the Schwarzschild metric. Paper~D37 (BH-1) established the
area law $\SPDL \propto \Rsurf$; D38 establishes the coefficient $4\pi$.
The ascending direction would close both results simultaneously.

\begin{openprob}[Axiom-level derivation of BH-2]
\label{op:counting}
Derive the identity
\begin{equation}
\ln\,\Omega_{\mathrm{surf}}(N) = 4\pi\left(\frac{\Meff(N)}{\MPl}\right)^2
\label{eq:op1}
\end{equation}
from the four \PDL\ axioms C1--C4, without invoking Schwarzschild geometry.
This requires: (a) a precise combinatorial definition of $\Omsurf(N)$ as a
function of the active surface $\Rsurf$ and the ensemble size $N$; and (b)
a proof that this count equals $4\pi(\sig(N)\cdot N\cdot m_p/\MPl)^2$.

This problem is conditional on OP1 of D36 (algebraic derivation of
$\kap = \Rsurf/\Rtot$ from C1--C4): once $\kap$ is proved from the axioms,
$\Rsurf$ is fixed combinatorially and the counting problem is well-posed.
\end{openprob}

\subsection{Structural interpretation}

The identity $\SBH/k_B = 4\pi(\Meff/\MPl)^2$ has a clean reading: the entropy
counts the number of Planck-mass units squared in the gravitationally active
mass, multiplied by the sphere solid angle. In \PDL\ terms, $\Meff/\MPl$ is
the ratio of two structural scales: the proton mass amplified by $\sig(N)\cdot N$
(the effective gravitational mass), and the Planck mass $\sqrt{\hb c/\GPDL}$
determined by the single-proton bridge formula. Both are combinatorial objects
fixed by the quintuplet $(24,28,930,10087,11017)$.

The fact that their ratio squared gives the Bekenstein--Hawking entropy, with
no free parameter, is the strongest numerical signature so far that the \PDL\
proton quintuplet encodes the full thermodynamics of gravitational surfaces.

\begin{conjecture}[BH-3: axiom-level Bekenstein--Hawking]
\label{conj:BH3}
The Bekenstein--Hawking entropy is a theorem of the four \PDL\ axioms: there
exists a combinatorial definition of $\Omsurf(N)$ derived from C1--C4 such that
\begin{equation}
k_B\ln\,\Omsurf(N) = 4\pi\,k_B\left(\frac{\sig(N)\cdot N\cdot m_p}{\MPl}\right)^2.
\label{eq:BH3}
\end{equation}
If true, Conjecture~BH-3 would constitute a derivation of the Bekenstein--Hawking
formula from combinatorial first principles, making the coefficient $1/4$ (in
Planck units) a prediction of the proton's internal architecture.
\end{conjecture}

% ══════════════════════════════════════════════════════════════════════════════
\section{Epistemic status summary}
\label{sec:epistemic}

\begin{table}[ht]
\centering
\caption{Epistemic status of the results of D38.}
\label{tab:epistemic}
\medskip
\begin{tabular}{>{\raggedright}p{6.5cm}p{3.2cm}p{3.5cm}}
\toprule
Statement & Status & Dependency \\
\midrule
$\Geff(N) = \sig(N)\cdot\GPDL$
  & Theorem under H1, H2 & D31/D36 \\[3pt]
$\SBH/k_B = 4\pi(\Meff/\MPl)^2$ (BH-2)
  & Algebraic theorem & Gate~3 + Schwarzschild \\[3pt]
$C_0 = 4\pi(m_p/\MPl)^2$ exact
  & Corollary, verified $<10^{-15}$
  & D38 Corollary~\ref{cor:C0} \\[3pt]
Entropy suppression $\sig(N)^2$ at epoch $N$
  & Prediction (Pred.~\ref{pred:entropy_epoch})
  & Gate~3 \\[3pt]
PBH threshold shift $+11.9\%$
  & Prediction (Pred.~\ref{pred:PBH})
  & Gate~3 + Hawking evaporation \\[3pt]
$S/E$ ratio $\propto\sig(N(z))$
  & Prediction (Pred.~\ref{pred:SoverE})
  & Gate~3 \\[3pt]
$\ln\Omsurf = 4\pi(\Meff/\MPl)^2$ from axioms
  & Open (OP1 D38)
  & OP1 D36 + combinatorics \\[3pt]
Bekenstein--Hawking as axiom theorem (BH-3)
  & Conjecture
  & OP1 D38 \\
\bottomrule
\end{tabular}
\end{table}

% ══════════════════════════════════════════════════════════════════════════════
\section*{Conclusion}

This paper has established one algebraic theorem and three falsifiable
predictions.

\emph{Proposition BH-2.} From Gate~3 alone ($\Geff(N) = \sig(N)\cdot\GPDL$),
the Bekenstein--Hawking entropy of a \PDL\ black hole satisfies
$\SBH/k_B = 4\pi(\Meff/\MPl)^2$, where $\Meff = \sig(N)\cdot N\cdot m_p$
is the gravitationally active mass. The identity is algebraically exact and
verified to $<10^{-15}$ over ten decades of mass. The universal constant
$C_0 = 4\pi(m_p/\MPl)^2 = 7.4221\times10^{-38}$ contains no free parameter.

\emph{Physical reading.} The Bekenstein--Hawking entropy counts the square
of the gravitationally active mass in Planck units. The $4\pi$ factor is
geometric (sphere solid angle). The \PDL\ modification is multiplicative:
$\SBH^\PDL = \sig(N)^2\cdot\SBH^\mathrm{GR}$, with $\sig(N) < 1$ at high
redshift.

\emph{Predictions.} Three falsifiable consequences: (i) $28.6\%$ entropy
suppression at the CMB epoch; (ii) $11.9\%$ upward shift in the primordial
black hole survival threshold, testable via Fermi-LAT; (iii) redshift-dependent
$S/E$ ratio in AGN populations.

\emph{Open problem.} The ascending derivation --- proving
$\ln\Omsurf = 4\pi(\Meff/\MPl)^2$ from C1--C4 without geometric input ---
is OP1 of D38 and constitutes Conjecture BH-3. Its resolution would make
the Bekenstein--Hawking formula a theorem of the four \PDL\ axioms.

% ── Bibliography ─────────────────────────────────────────────────────────────
\begin{thebibliography}{99}

\bibitem{PDL_Main_2026}
C.~Laubscher,
\textit{The Projective Dynamic Logo: A Structural Derivation of Fundamental Constants
from Finite Signed Graphs},
Zenodo, 2026. \url{https://doi.org/10.5281/zenodo.14999037}

\bibitem{PDL_D25_2026}
C.~Laubscher,
\textit{The Surface Coverage Fraction $\sigma(N)$ and the Effective Gravitational
Coupling in the PDL Framework} (D25), Zenodo, 2026.

\bibitem{PDL_D31_2026}
C.~Laubscher,
\textit{Proof of the Density-Dependent Gravitational Coupling
$G_{\mathrm{eff}}(N)=\sigma(N)\cdot G_{\mathrm{PDL}}$ in the PDL Framework} (D31),
Zenodo, 2026. \url{https://doi.org/10.5281/zenodo.19294984}

\bibitem{PDL_D36_2026}
C.~Laubscher,
\textit{Proof of $G_{\mathrm{eff}}(N)=\sigma(N)\cdot G_{\mathrm{PDL}}$ from the
Trace Structure of $K_4$ and Independence of Gravitational Engagement} (D36),
Zenodo, 2026. \url{https://doi.org/10.5281/zenodo.19323033}

\bibitem{PDL_D37_2026}
C.~Laubscher,
\textit{Surface Locality of Relational Entropy and the PDL Area Law} (D37),
Zenodo, 2026.

\bibitem{Bekenstein_1973}
J.~D.~Bekenstein,
\textit{Black holes and entropy},
Phys.\ Rev.\ D \textbf{7}, 2333 (1973).
\url{https://doi.org/10.1103/PhysRevD.7.2333}

\bibitem{Hawking_1975}
S.~W.~Hawking,
\textit{Particle creation by black holes},
Commun.\ Math.\ Phys.\ \textbf{43}, 199 (1975).
\url{https://doi.org/10.1007/BF02345020}

\bibitem{CODATA2022}
P.~J.~Mohr, D.~B.~Newell, B.~N.~Taylor, E.~Tiesinga,
\textit{CODATA recommended values of the fundamental physical constants: 2022},
Rev.\ Mod.\ Phys.\ \textbf{96}, 025002 (2024).

\bibitem{FermiLAT}
M.~Ackermann et al.\ (Fermi-LAT Collaboration),
\textit{Searching for Dark Matter Annihilation from Milky Way Dwarf Spheroidal
Galaxies with Six Years of Fermi Large Area Telescope Data},
Phys.\ Rev.\ Lett.\ \textbf{115}, 231301 (2015).

\end{thebibliography}

\end{document}